% ============================================================================
%  บทที่ 6: สรุปผลและข้อเสนอแนะ
% ============================================================================

\chapter{สรุปผลและข้อเสนอแนะ}
\label{ch:conclusion}

% ============================================================================
\section{สรุปผลการดำเนินงาน}
\label{sec:conclusion}
% ============================================================================

โครงงานนี้ได้ออกแบบและพัฒนาหุ่นยนต์หยิบและวางชิ้นงานแบบวงกลม (Circular Pick and Place Robot)
สำหรับลำเลียงชิ้นงาน Connection Rod ในกระบวนการผลิตเชิงอุตสาหกรรม
โดยครอบคลุมวิศวกรรมทุกมิติของระบบ Cyber-Physical

\textbf{ด้านเชิงกล}
ระบบส่งกำลังแบบ 2 ชั้น (Planetary Gearbox $N_g = 24$ + Belt Drive HTD5M $N_b = 2.917$)
ให้อัตราทดรวม $N = 70$ และประสิทธิภาพ $\eta = 0.836$
การออกแบบเพลาตาม ASME MSST ได้ขนาดขั้นต่ำ 10.59\,mm
และเลือกใช้ 17\,mm ที่จุดวิกฤต (FOS $\geq 1.6$)
การวิเคราะห์ FEA พบการโก่งตัวสูงสุด 0.2281\,mm และ FOS ต่ำสุด 11.03
ซึ่งเป็นไปตามข้อกำหนด M3 ทุกประการ

\textbf{ด้านไฟฟ้า}
ตู้ควบคุมออกแบบตามมาตรฐานอุตสาหกรรม
แบ่ง Power Bus ออกเป็น High Power (24\,V) และ Logic (5\,V) อย่างชัดเจน
E-Stop ทำงานในระดับ Hardware แบบ Fail-Safe
และมีวงจรป้องกัน Overcurrent, Surge และ Short Circuit ตามข้อกำหนด E6 และ E7

\textbf{ด้านการควบคุม}
การระบุระบบด้วย Locked Rotor Test และ Chirp Excitation
ให้ค่าพารามิเตอร์ที่สมบูรณ์ครบ 8 ตัว
ตัวควบคุม Cascade PID ออกแบบด้วย Root Locus
ได้ Settling Time 0.450\,s และ Overshoot 0.00\% จากการ Simulation
ZV Input Shaper ลด Rod Settling Time จาก 6.44\,s เหลือ $\approx 0$\,s ต่อ Move
ทำให้ Cycle Time รวมเหลือ 30\,s (Simulation) ซึ่งเป็นไปตาม P1
Kalman Filter 4 สถานะใช้ Steady-State Gain จาก DARE
ลด Computational Load บน STM32 ได้อย่างมีประสิทธิภาพ

\textbf{ด้านซอฟต์แวร์และการบูรณาการ}
Firmware บน STM32G474RE ประกอบด้วยโมดูลครบถ้วน
ได้แก่ Cascade PID, ZV Shaper, Kalman Filter, FSM, Modbus RTU
และ Joystick Interface
การสื่อสารกับ Base System ผ่าน Modbus RTU (19200 bps, 8E1, Slave Address~21)
พร้อม Python asyncio WebSocket Backend

% ============================================================================
\section{ปัญหา อุปสรรค และบทเรียนที่ได้}
\label{sec:lessons_learned}
% ============================================================================

\subsection{ปัญหาเชิงกล}
\label{subsec:mech_lessons}

\begin{itemize}
  \item \textbf{ความคลาดเคลื่อนของชิ้นส่วนที่ผลิตภายนอก:}
        รูไม่ตรงแนวศูนย์และผิวหน้าไม่ได้รับการปาดตามแบบ
        สะท้อนให้เห็นความสำคัญของการ Inspection
        ชิ้นส่วนก่อนประกอบ และการสื่อสาร Tolerance ให้ชัดเจนกับผู้ผลิต

  \item \textbf{Laser Cut Tolerance:}
        ชิ้นส่วนที่ตัดด้วย Laser Cut มี Tolerance สะสมที่ส่งผลต่อการ Align
        โดยเฉพาะที่จุดหมุนของแขนกล
        แนะนำให้กำหนด Fitment Tolerance อย่างชัดเจนในแบบ
        และทำ Shimming ระหว่างการประกอบ

  \item \textbf{ขนาด Proximity Sensor ไม่ตรงกับแบบ:}
        เกิดจากการเลือก Sensor รุ่นอื่นในภายหลัง
        โดยไม่ได้อัพเดตแบบ CAD ก่อนส่งผลิต
        ควรล็อก BOM ก่อน Finalize แบบ
\end{itemize}

\subsection{ปัญหาด้านซอฟต์แวร์และการบูรณาการ}
\label{subsec:sw_lessons}

\begin{itemize}
  \item \textbf{Abnormal Heartbeat ใน Demo~1:}
        Backend \texttt{main.exe} ของ Base System มีปัญหาคอขวด
        ทำให้ Heartbeat Timeout และระบบ Handshake ไม่สำเร็จ
        แก้ไขโดยเปลี่ยน Backend เป็น Python asyncio WebSocket
        บทเรียน: ควรทดสอบ End-to-End Integration
        กับ Base System จริงก่อน Demo อย่างน้อย 1 สัปดาห์

  \item \textbf{Joystick Mode Blocking Bug:}
        พบว่าคำสั่ง Joystick บางตัวไม่ถูก Block
        เมื่อระบบอยู่ใน Base System Mode
        เนื่องจาก Mode Check Logic ไม่ครอบคลุมทุก Command Path
        แก้ไขโดยเพิ่ม Guard Condition ที่ Entry Point ของ
        \texttt{Motor\_ProcessCommand}

  \item \textbf{Gripper พังระหว่าง Demo~2:}
        Base System ส่งคำสั่งตำแหน่งผิดพลาด
        ทำให้แขนกลเคลื่อนที่ผิดทิศและชน Connection Rod
        บทเรียน: ควรมี Software Workspace Limit
        ที่ตรวจสอบทุก Incoming Command จาก Base System
        ก่อน Execute เสมอ
        และควรทดสอบ Safety Boundary ให้ครบก่อน Full Demo
\end{itemize}

\subsection{ปัญหาด้านการออกแบบระบบควบคุม}
\label{subsec:ctrl_lessons}

\begin{itemize}
  \item \textbf{SysID Preprocessing vs Raw Data:}
        การกรองสัญญาณด้วย Butterworth $f_c = 10$\,Hz ก่อน Estimation
        ให้ผลดีกว่าข้อมูลดิบ แต่มีความเสี่ยงจาก Phase Shift
        ที่แก้ไขด้วยการใช้ \texttt{filtfilt}
        การตัดสินใจนี้ถูกต้องสำหรับโครงงานนี้
        แต่ควรรายงาน Validation ให้ชัดเจนในทุกกรณีที่ Preprocess ข้อมูล

  \item \textbf{Bandwidth Separation Heuristic:}
        การตั้ง Inner Loop เร็วกว่า Outer Loop 6 เท่าเป็น Rule of Thumb
        ไม่ใช่เกณฑ์ที่เข้มงวด
        ในการออกแบบที่ต้องการความแม่นยำสูงขึ้น ควรออกแบบ Inner Loop ก่อน
        แล้ว Model Closed-loop Inner Loop เป็น Plant ที่แท้จริง
        สำหรับการออกแบบ Outer Loop ต่อไป

  \item \textbf{ZV Shaper Sensitivity:}
        ZV Shaper มีความไวต่อความคลาดเคลื่อนของ $\omega_{n,rod}$
        หากค่าจริงต่างจากที่ออกแบบ $\pm 5\%$
        Residual Vibration จะเพิ่มขึ้น
        สำหรับระบบที่ต้องการความทนทานสูง
        ควรพิจารณาใช้ ZVD (3-Impulse) แทน ZV (2-Impulse)
\end{itemize}

% ============================================================================
\section{ข้อเสนอแนะสำหรับการพัฒนาต่อยอด}
\label{sec:future_work}
% ============================================================================

\subsection{การพัฒนาระบบควบคุม}
\label{subsec:future_control}

\begin{itemize}
  \item \textbf{Hardware Tuning ของ PID Gains:}
        ค่า Gain ที่ใช้งานปัจจุบันมาจากการออกแบบบน Reduced-Order Model
        ควรทำ Hardware Commissioning เพื่อปรับ Gain Fine-tuning
        โดยเฉพาะ $K_{i,vel}$ ที่มีขนาดใหญ่
        และตรวจสอบ Anti-windup ว่าทำงานถูกต้องภายใต้ Load จริง

  \item \textbf{เปิดใช้งาน Feedforward Gains:}
        $K_{vff} = 3.033$\,V/(rad/s) และ $K_{aff} = 0.445$\,V/(rad/s$^2$)
        ถูก Disable ระหว่าง Commissioning
        เมื่อ SysID Uncertainty ต่ำกว่า 10\% แล้ว
        ควรเปิดใช้งานเพื่อลด Tracking Error ระหว่างการเคลื่อนที่

  \item \textbf{Upgrade เป็น ZVD Shaper:}
        ZVD (3-Impulse) มีความทนทานต่อความคลาดเคลื่อนของ $\omega_n$
        มากกว่า ZV ควรพิจารณาเมื่อทราบ Uncertainty ของ $\omega_{n,rod}$
        จากการวัดจริงแล้วว่าเกิน 10\%

  \item \textbf{Online Parameter Estimation:}
        ในอนาคตสามารถ Implement Recursive Least Squares (RLS)
        เพื่อ Update $\omega_{n,rod}$ แบบ Online
        เมื่อสภาพของ Connection Rod เปลี่ยนแปลงตามการใช้งาน
\end{itemize}

\subsection{การพัฒนาระบบกล}
\label{subsec:future_mech}

\begin{itemize}
  \item \textbf{ลด Backlash ของ Belt Drive:}
        ออกแบบ Tensioner ที่ปรับระยะได้ (Adjustable Idler)
        เพื่อให้สามารถรักษา Belt Tension ได้อย่างสม่ำเสมอ
        ตลอดอายุการใช้งาน

  \item \textbf{ปรับปรุงความแข็งแกร่งของ Gripper Mounting:}
        เพิ่ม Reinforcement ที่จุดยึด Gripper
        เพื่อป้องกันการเสียหายจากการชนในกรณี Software Fault
\end{itemize}

\subsection{การพัฒนาระบบไฟฟ้าและซอฟต์แวร์}
\label{subsec:future_elec_sw}

\begin{itemize}
  \item \textbf{Implement Software Workspace Limit:}
        เพิ่มการตรวจสอบ Incoming Command ทุกตัวจาก Base System
        ว่าอยู่ภายใน Safe Workspace ก่อน Execute
        เพื่อป้องกันการชนซ้ำ

  \item \textbf{Upgrade การสื่อสาร Gripper:}
        ในอนาคตสามารถ Implement CAN Bus (FDCAN1)
        เพื่อให้รองรับ Gripper Node ที่ซับซ้อนขึ้น
        พร้อม Watchdog และ Feedback ที่ครบถ้วนกว่า Digital I/O

  \item \textbf{Data Logging และ Monitoring:}
        เพิ่ม SD Card หรือ Wireless Telemetry
        สำหรับบันทึก Log ระยะยาว
        เพื่อ Predictive Maintenance และวิเคราะห์ Drift ของ Bearing
\end{itemize}